\section{Композиция}
\label{sec:composition}

Архитектура \texttt{PySATL Core} организована вокруг набора относительно независимых компонент, объединённых общей вычислительной моделью: характеристики распределения вычисляются через единые контракты, а правила вывода оформляются как расширяемый граф зависимостей. Такая композиция позволяет наращивать систему (новыми семействами, характеристиками, способами вычислений и сэмплирования) без модификации уже реализованных частей.

\subsection{Состав компонент}

На уровне верхней декомпозиции выделяются следующие компоненты (рис.~\ref{fig:components}):
\begin{itemize}[compact]
  \item \textbf{\texttt{families}} --- описание параметрических семейств распределений, их параметризаций и ограничений, а также фабрика построения конкретных распределений;
  \item \textbf{\texttt{distributions}} --- единая «поверхность» доступа к распределению: вычисление характеристик, выбор стратегий, работа с поддержкой (support) и интеграция с реестром характеристик;
  \item \textbf{\texttt{sampling}} --- инфраструктура генерации выборок и представления результата (тип выборки), поставляемая как отдельная, переиспользуемая часть ядра;
  \item \textbf{\texttt{transformations}} --- компонент описания преобразований распределений и правил пересчёта характеристик при композиции (на момент актуализации проектно значим и рассматривается как целевое расширение).
\end{itemize}

Внутри компонентов выделены подсистемы, фиксирующие устойчивые точки расширения:
\begin{itemize}[compact]
  \item в \texttt{distributions} --- \texttt{registry}, содержащий глобальный реестр преобразований характеристик и механизмы получения \emph{локального} представления графа для конкретного распределения;
  \item в \texttt{families} --- \texttt{builtins}, включающий готовые семейства и общие элементы описания параметризаций.
\end{itemize}

\begin{figure}[H]
  \centering
  \includegraphics[width=.65\linewidth]{components.png}
  \caption{Компоненты ядра и границы публичных API}
  \label{fig:components}
\end{figure}

\subsection{Роли компонент}

\paragraph{\texttt{families}.}
Компонент отвечает за \emph{порождение} распределений. Семейство задаётся как объект, содержащий набор параметризаций и правила их приведения к базовой форме. Для каждой параметризации могут быть объявлены ограничения на параметры. Результатом работы \texttt{families} является объект распределения, который реализует контракт \texttt{Distribution} и тем самым бесшовно подключается к вычислительной модели \texttt{distributions}.

\paragraph{\texttt{distributions}.}
Компонент определяет контракт распределения и обеспечивает вычисление характеристик. В частности, он фиксирует:
\begin{itemize}[compact]
  \item \emph{идентификаторы характеристик} (на уровне типов) и единый механизм запроса вычисления;
  \item \emph{стратегии} вычислений и сэмплирования как взаимозаменяемые политики;
  \item \emph{поддержку} случайной величины как отдельный объект со своими инвариантами;
  \item \emph{реестр характеристик} как расширяемый ориентированный граф правил вывода.
\end{itemize}
Компонент не фиксирует конкретный численный алгоритм «внутри» каждого метода, но задаёт формат его подключения и условия применимости.

\paragraph{\texttt{sampling}.}
Компонент инкапсулирует генерацию выборок и представление результата. Конкретные алгоритмы сэмплирования поставляются через \texttt{SamplingStrategy}; при этом распределение хранит выбранную стратегию как часть своего состояния и делегирует ей операцию \texttt{sample}.

\paragraph{\texttt{transformations}.}
Компонент предназначен для описания преобразований распределений (например, аффинные преобразования, смеси, свёртки) и их влияния на характеристики и поддержку. В целевой архитектуре такие правила должны быть выражены через те же абстракции вычислительных методов, что и «внутренние» преобразования характеристик, и регистрироваться в общем графе. Это обеспечивает единый путь вычисления характеристик как для атомарных распределений, так и для составных.

\subsection{Потоки данных и точки соприкосновения}

Типовой поток использования выглядит следующим образом:
\begin{enumerate}[compact]
  \item Пользователь выбирает семейство и задаёт параметры в удобной параметризации.
  \item \texttt{families} валидирует параметры и приводит их к базовой параметризации.
  \item Создаётся объект распределения, предоставляющий \texttt{Distribution}-API.
  \item При запросе характеристики \texttt{distributions} выбирает подходящий метод: либо из аналитически заданных вычислений, либо через вывод по графу реестра; при необходимости применяется кэширование.
  \item При запросе выборки распределение делегирует генерацию выборки стратегии \texttt{SamplingStrategy}.
\end{enumerate}

Таким образом, граница между компонентами проводится по \emph{контрактам}: \texttt{families} формирует корректный экземпляр распределения, а \texttt{distributions} обеспечивает вычислительную инфраструктуру и расширяемость через реестр и стратегии.
